\documentclass[11pt]{article}

\usepackage[margin=1in]{geometry}
\usepackage{amsmath, amssymb, amsthm}
\usepackage{bm}
\usepackage{booktabs}
\usepackage{algorithm}
\usepackage{algpseudocode}
\usepackage{xcolor}
\usepackage[colorlinks=true, linkcolor=blue, citecolor=blue]{hyperref}

% ---- notation shorthands -------------------------------------------------
\newcommand{\R}{\mathbb{R}}
\newcommand{\E}{\mathbb{E}}
\newcommand{\st}{s}
\newcommand{\Vnorm}{\widetilde{V}}
\newcommand{\theem}{\bm{\theta}}
\newcommand{\skill}{z}
\newcommand{\given}{\,\vert\,}
\newcommand{\argmax}{\operatorname*{arg\,max}}
\newcommand{\softmax}{\operatorname{softmax}}
\newcommand{\Bern}{\mathrm{Bernoulli}}
\newcommand{\Unif}{\mathcal{U}}
\newcommand{\code}[1]{\texttt{\small #1}}

\title{\textbf{An Adaptive Skill-Switching Transformer:\\
HG-DAgger inside a Bayes-Adaptive MDP of Synthetic Failures}}
\author{SMERL project --- technical note}
\date{\today}

\begin{document}
\maketitle

\begin{abstract}
We formalize a method for training a transformer policy that \emph{adapts on-line}:
it commits to a behavioral skill, monitors a scalar value signal, and abandons the
skill for a better one when that signal stalls. The skill repertoire is a fixed set
of $K$ skill-conditioned controllers (e.g.\ a SMERL actor on a $2$D point-goal task).
Adaptation is learned, not hand-coded, by three coupled pieces. (i)~A
\emph{Bayes-Adaptive MDP} (BAMDP) injects \emph{synthetic} failures: a latent
per-skill failure rate $\theem$ is drawn each meta-episode, and a hazard process
corrupts the exposed value so a doomed run is visible as a value that stops climbing.
(ii)~An \emph{HG-DAgger} loop collects on-policy intervention demonstrations in which
an expert takes over with a good skill once the value has stalled. (iii)~A causal
\emph{multimodal transformer} is trained by behavior cloning on the merged data with
a sequence of state, skill, and value tokens, predicting the skill from the state and
the action from the skill. At inference the model re-reads its own skill head each
step; a change in the committed skill is a switch. We give the exact generative
process, loss, and rollout procedure as implemented.
\end{abstract}

\tableofcontents

% =========================================================================
\section{Setup and notation}
\label{sec:setup}

\paragraph{Base MDP.}
The base environment is a deterministic $2$D point-mass goal-reaching MDP with state
$\st=(x,y,\dot x,\dot y,g_x,g_y)\in\R^{6}$, continuous action $a\in[-1,1]^2$
(acceleration), fixed start and goal, and a sparse-ish reward
$r(\st,a)=-\lVert (x,y)-(g_x,g_y)\rVert - c\lVert a\rVert^2$ with a success bonus when
the agent enters the goal radius. Episodes have horizon $H$ (we use $H=100$).

\paragraph{Skills.}
A pretrained SMERL actor provides $K$ skill-conditioned controllers
$\pi(a\given \st,\skill)$, $\skill\in\{0,\dots,K-1\}$. Each skill is a distinct,
near-deterministic route to the goal; we write $\pi(\st,\skill)$ for its deterministic
action $\argmax$. The skills are \emph{given} and frozen throughout.

\paragraph{Skill-conditioned value.}
A skill-conditioned state value $V(\st,\skill)\in\R$ is trained by Monte-Carlo
regression on success-returns (\code{train\_task\_value.py}). It estimates how well
skill $\skill$ is doing from state $\st$; it climbs toward its maximum as the skill
nears the goal. We use a skill-conditioned $V(\st,\skill)$ rather than a single
task value $V(\st)$ because the former gives a much steeper, skill-identifiable
climb (the chosen ablation).

\paragraph{Per-skill value normalization.}
To make a stall maximally legible we rescale each skill's value to $[0,1]$:
\begin{equation}
\Vnorm(\st,\skill) \;=\;
\operatorname{clip}\!\left(\frac{V(\st,\skill)-v^{\min}_{\skill}}
{v^{\max}_{\skill}-v^{\min}_{\skill}},\, 0,\, 1\right),
\label{eq:vnorm}
\end{equation}
where $(v^{\min}_{\skill},v^{\max}_{\skill})$ are the $[1,99]$ percentiles of
$V(\cdot,\skill)$ over deterministic rollouts of skill $\skill$ from perturbed starts
(\code{compute\_skill\_value\_norm}). This steepens the climb so a frozen value is an
unambiguous plateau. Skills not in the normalized set pass through unchanged.

% =========================================================================
\section{The Bayes-Adaptive MDP of synthetic failures}
\label{sec:bamdp}

The BAMDP (\code{bamdp\_env.py}, class \code{SyntheticFailureBAMDP}) wraps the base
env and injects failures whose statistics depend on a hidden latent. It is the
training ground in which a stall \emph{means} ``this skill will not succeed.''

\subsection{Latent and meta-episodes}
A \emph{meta-episode} fixes a latent vector of per-skill \emph{failure rates}
\begin{equation}
\theem=(p_1,\dots,p_K), \qquad p_i \sim \mathrm{Beta}(\alpha,\beta)\ \text{i.i.d.}
\end{equation}
($\alpha=\beta=0.5$ by default). In the two-skill regime we actually train on, we
override the prior with a structured draw: one of the in-play skills is made
\emph{bad}, $p_{\mathrm{bad}}=0.95$, the rest near-safe, $p_i=0.02$. The latent
$\theem$ is what a Bayes-optimal agent must infer from the value stream; it is not
observed. Each failure rate is converted to a total \emph{hazard budget}
\begin{equation}
H_i \;=\; -\ln(1-p_i),
\end{equation}
so that spending the whole budget over an episode reproduces the marginal failure
probability $p_i$.

\subsection{Hazard injection (the \texttt{certain\_skill} mechanism)}
At each step the wrapper forms a belief $w\in\Delta^{K-1}$ over which skill is
currently executing. In general $w=\softmax(\text{Disc}(\st))$ from a behavioral
discriminator; in the simplified regime we use (\code{certain\_skill}), the running
skill is \emph{known}, so $w$ is a one-hot on the active skill $\skill_{\mathrm{act}}$:
\begin{equation}
w \;=\; e_{\skill_{\mathrm{act}}}.
\end{equation}
Under the \emph{budget} schedule, the per-skill instantaneous hazard spends the
remaining budget over the estimated remaining steps,
\begin{equation}
h_i \;=\; \frac{H^{\mathrm{rem}}_i}{\max(r^{\mathrm{rem}}_i,\, 1)},\qquad
h \;=\; \textstyle\sum_i w_i\, h_i,\qquad
f \;=\; 1-e^{-h},
\end{equation}
and a forced failure is drawn $F_t\sim\Bern(f)$. The budget is then debited in
proportion to presence, $H^{\mathrm{rem}}_i \leftarrow \max(H^{\mathrm{rem}}_i - w_i
h_i,0)$ and $r^{\mathrm{rem}}_i \leftarrow r^{\mathrm{rem}}_i - w_i$. Because $w$ is a
one-hot, the active skill draws failures at exactly its own hazard $h_{\skill_{\mathrm
{act}}}$ --- the discriminator is treated as certain ($p=1$).

\subsection{Absorbing failure and value corruption}
Failures are \emph{continuing} (\code{continue\_on\_failure}): a forced failure does
\emph{not} end the episode. Instead it is an absorbing event that decouples the
\emph{exposed} value $v^{\mathrm{obs}}_t$ from the true state value. While healthy,
\begin{equation}
v^{\mathrm{obs}}_t \;=\; \Vnorm(\st_t,\skill_{\mathrm{act}}).
\end{equation}
At the failure instant $t_f$ we freeze $v^\star=\Vnorm(\st_{t_f},\skill_{\mathrm{act}})$
and pick a corruption mode $\in\{\text{plateau},\text{decline}\}$ with probability
$\tfrac12$ each; thereafter
\begin{equation}
v^{\mathrm{obs}}_{t_f+k} \;=\;
\begin{cases}
v^\star & \text{plateau},\\[2pt]
v^\star\,\gamma^{\,k} & \text{decline}\ (\gamma=0.95).
\end{cases}
\label{eq:corrupt}
\end{equation}
A doomed run is therefore visible as a value that stops climbing toward $1$ (plateau)
or sags (decline). Once failed, no further failures are injected (the event is
absorbing) and the goal bonus is suppressed, so the run truncates at $H$ without
success.

\subsection{Observation interface}
With value exposure on, the wrapper returns a \emph{dict} so each consumer reads the
right stream:
\begin{equation}
o_t=\big\{\,\code{state}=\st_t\in\R^{6},\ \ \code{value}=v^{\mathrm{obs}}_t\in[0,1],
\ \ \code{failing}=\mathbf{1}[\text{failed}]\in\{0,1\}\,\big\}.
\end{equation}
The skill controller is driven with \code{state} only (it was trained on the bare
$6$-D observation); \code{value} is the adaptation signal; \code{failing} is used only
for bookkeeping during data collection, never at test time.

\subsection{Switching and intervention semantics}
Two operations change the active skill mid-episode.
\begin{itemize}
\item \textbf{Agent switch} (test time, \code{switch\_skill}): re-evaluate failure
under the new skill --- clear the absorbing flag and refresh $v^{\mathrm{obs}}$ ---
but \emph{keep injecting} failures, because the new skill may also be bad.
\item \textbf{Expert intervention} (training time, \code{intervene}): same, but also
\emph{stop} injecting failures for the rest of the sub-episode, so the good skill runs
cleanly to the goal.
\end{itemize}
In both cases, with \code{reset\_on\_switch} enabled the environment \emph{teleports}
to a fresh start (a new draw of the base env) and resets the per-attempt hazard/value
state, while \emph{keeping the meta step clock} $t$. Operationally, a switch is
``give up and retry from scratch with the other skill,'' the two attempts sharing one
horizon budget $H$. This makes a switch a genuine state \emph{and} value
discontinuity, which the transformer must key on.

% =========================================================================
\section{HG-DAgger data collection}
\label{sec:dagger}

We collect two stores and merge them. A \emph{base} store of clean
skill demonstrations (with DART action noise on the environment) anchors the action
and skill heads; a \emph{DAgger} store of on-policy interventions teaches the switch.

\subsection{DART base demonstrations}
For each skill we roll the deterministic teacher $\pi(\st,\skill)$ and store the clean
action as the target, while the \emph{environment} is driven by a DART-noised action
\begin{equation}
\tilde a_t \;=\; \operatorname{clip}\big(\pi(\st_t,\skill) + \epsilon_t,\,-1,1\big),
\qquad \epsilon_t\sim\Unif(-c,c)^2,
\end{equation}
with $c$ a small fraction (5\%) of the mean action magnitude. Recording the clean
action under noised states is the DART recipe: it widens state coverage so the cloned
policy is robust to its own drift, without biasing the action targets.

\subsection{On-policy interventions}
Each DAgger episode (Algorithm~\ref{alg:dagger}) is an HG-DAgger rollout: the
\emph{transformer under training} chooses the skill to attempt, the teacher executes
it, the BAMDP decides whether it fails, and an \emph{expert} intervenes once the
failure has been visible for a few steps.

\begin{algorithm}[t]
\caption{Collect one HG-DAgger intervention episode (\code{collect\_dagger\_episode})}
\label{alg:dagger}
\begin{algorithmic}[1]
\Require BAMDP with latent $\theem$ (one bad skill); transformer $M$; teacher $\pi$;
expert good skill $\skill_{\mathrm{good}}$; window $[k_{\min},k_{\max}]$; delay $d$.
\State $o\gets\textsc{BamdpReset}()$; \ $\st_0\gets o.\code{state}$
\State $\skill_0\sim \softmax\big(M_{\mathrm{skill}}(\st_0)\big)$
\Comment{on-policy: attempt the model's own skill}
\If{$\skill_{\mathrm{good}}=\skill_0$} re-draw $\skill_{\mathrm{good}}$ from the top-2
of $1-\theem$, $\neq \skill_0$ \EndIf
\State $\textsc{SetActiveSkill}(\skill_0)$; \ active $\gets\skill_0$;\ intervened
$\gets$ false;\ $K\sim\Unif\{k_{\min},\dots,k_{\max}\}$
\While{episode not done}
  \State $a_t \gets \pi(\st_t,\text{active})$
  \Comment{clean action = \emph{target}}
  \State $\tilde a_t \gets \operatorname{clip}(a_t+\epsilon_t,-1,1)$;\quad
         $o',r,\text{done}\gets\textsc{BamdpStep}(\tilde a_t)$
  \State record $a_t$, active skill, $v^{\mathrm{obs}}_t$, \code{failing}$_t$
  \If{$\neg$intervened \textbf{and} \code{failing}$_t$}
     \State $k_{\mathrm{fail}}\!\mathrel{+}=\!1$;\quad
            \textbf{if } $k_{\mathrm{fail}}\ge K$ \textbf{ then }
            $\textsc{Intervene}(\skill_{\mathrm{good}})$; active $\gets
            \skill_{\mathrm{good}}$; intervened $\gets$ true
  \EndIf
  \State $o \gets \textsc{Observe}()$ \textbf{if} just-switched (re-read the fresh
         teleported start) \textbf{else} $o'$
\EndWhile
\State \textbf{keep} the episode iff (intervened \textbf{and} success)
\State \Return episode dict with per-step \code{skills} and decoupled
\code{skill\_target}
\end{algorithmic}
\end{algorithm}

\paragraph{Choosing the takeover skill.}
The expert's good skill is sampled from the two highest success rates
$1-\theem$ (so it is among the genuinely-good options), forced to differ from
$\skill_0$ so a real switch occurs (\code{pick\_good\_skill}).

\paragraph{Decoupled skill target.}
This is the key labeling trick. The \emph{executed} skill token in the sequence stays
$\skill_0$ until the physical switch, but the \emph{target} skill label flips to
$\skill_{\mathrm{good}}$ a few steps after the failure becomes observable:
\begin{equation}
\code{skill\_target}_t =
\begin{cases}
\skill_0 & t < t_f-1+d,\\
\skill_{\mathrm{good}} & t \ge t_f-1+d,
\end{cases}
\qquad d=2\ \text{(observation delay)}.
\end{equation}
So while the value is still stalling (executed skill $\skill_0$ remains in context),
the supervised \emph{decision} is already ``switch to $\skill_{\mathrm{good}}$.'' This
decoupling is what teaches the skill head to act on the stall rather than merely to
repeat whatever skill produced the action. (The ablation that couples target to the
executed skill never learns to switch.)

\paragraph{Keep filter and merge.}
Only episodes that intervened \emph{and} reached the goal are kept --- the dataset
shows the policy a stall followed by a \emph{successful} recovery. The kept DAgger
store is concatenated with the DART base store; \code{SeqDataset} interleaves episodes
from both directories.

% =========================================================================
\section{The multimodal transformer and its training}
\label{sec:transformer}

\subsection{Token layout}
Each timestep is expanded into three tokens, in the order \emph{state, skill, value}:
\begin{equation}
\underbrace{\st_0}_{\text{state}}\ \underbrace{\skill_0}_{\text{skill}}\
\underbrace{v_0}_{\text{value}}\ \ \st_1\ \skill_1\ v_1\ \ \cdots
\end{equation}
The supervision pattern (\code{pattern: [[state, skill], [skill, action], value]}) is:
\begin{itemize}
\item from the \textbf{state} token, predict the \textbf{skill} $\skill_t$ (this is
the switch decision);
\item from the \textbf{skill} token, predict the \textbf{action} $a_t$;
\item the \textbf{value} token is an input only (never predicted).
\end{itemize}
Placing the skill token \emph{before} the value and action makes the policy
\emph{express} which skill it is running, and lets the action head condition on it.

\subsection{Backbone}
A self-contained, minGPT-style causal transformer (\code{TrajectoryGPT}) consumes
\emph{continuous} token embeddings: one linear (or, for the discrete skill, an
embedding-lookup) encoder per modality maps a raw token to $\R^{E}$, a learned
\emph{type} embedding per modality is added, and learned positional embeddings and
causal masking complete the input. We use $E=256$, $6$ layers, $4$ heads, dropout
$0.1$. The hidden at position $i$ predicts token $i{+}1$ through \emph{that target's}
modality head, with a per-position loss mask selecting which predictions count.

\subsection{Heads and per-modality losses}
\paragraph{Skill head (categorical).}
A linear map to $K$ logits; cross-entropy against the (decoupled) target skill:
\begin{equation}
\ell^{\mathrm{skill}}_t \;=\; \mathrm{CE}\big(M_{\mathrm{skill}}(\st_t),\,
\code{skill\_target}_t\big).
\end{equation}

\paragraph{Action head (discretized, ``disc'').}
Rather than a Gaussian (whose NLL gradients are pathological here), each action
dimension is a categorical over $B$ uniform bins on $[-1,1]$ ($B=51$). For target
$a_t\in[-1,1]^{2}$ the bin index is
$\iota_{t,j}=\big\lfloor \tfrac{(a_{t,j}+1)}{2}B\big\rfloor$ (clamped), and
\begin{equation}
\ell^{\mathrm{act}}_t \;=\; \sum_{j=1}^{2}
\mathrm{CE}\big(M_{\mathrm{act}}(\skill_t)_{j},\,\iota_{t,j}\big).
\end{equation}
At rollout the action is read back as the bin-center of the $\argmax$ per dimension.

\paragraph{Two layers of class balancing.}
The switch decision is rare, so we upweight it at two levels.
\begin{itemize}
\item \emph{Layer 1 (per-position).} A switch-decision position --- a state token whose
executed skill changes next step, $\skill_{t+1}\neq\skill_t$ --- gets loss weight
$\omega=20$; all other positions weight $1$. Within a modality the loss is the
weighted mean
\begin{equation}
\mathcal{L}_{\text{mod}} \;=\;
\frac{\sum_{t\in\text{mod}} \omega_t\,\ell_t}{\sum_{t\in\text{mod}} \omega_t}.
\end{equation}
\item \emph{Layer 2 (per-modality).} An optional multiplier $\lambda_{\text{mod}}$
keeps the action NLL from dominating the shared trunk; the token loss is
\begin{equation}
\mathcal{L}_{\mathrm{tok}} \;=\; \sum_{\text{mod}\in\{\text{skill},\text{act}\}}
\lambda_{\text{mod}}\,\mathcal{L}_{\text{mod}}.
\end{equation}
\end{itemize}

\subsection{The logic gate (a training-dynamics probe)}
To verify that the \emph{hidden representation} at the state token actually contains
enough signal to detect a switch, we attach an auxiliary binary head
$g=\sigma(M_{\mathrm{logic}}(h_{\mathrm{state}}))$ predicting ``does a switch happen
here?'' It is trained with weighted BCE \emph{restricted to switch-containing
episodes} (the pure no-switch DART episodes are all-negative and would bias it to
``never switch''), with positive class weight $\rho=50$:
\begin{equation}
\mathcal{L}_{\mathrm{logic}} \;=\;
\mathrm{BCE}_{\rho}\!\big(M_{\mathrm{logic}}(h_{\mathrm{state}}),\,
\mathbf{1}[\text{switch}]\big),
\qquad
\mathcal{L} \;=\; \mathcal{L}_{\mathrm{tok}} + \lambda_{\mathrm{logic}}\,
\mathcal{L}_{\mathrm{logic}}.
\end{equation}
Crucially the logic gate is a \emph{diagnostic only}: it shares the backbone so its
high recall ($\approx0.9$) certifies the switch is decodable, but it is \emph{never}
consulted at inference. Comparing runs, when the logic gate confirms the signal is
present while the skill head still won't switch, the cure is the decoupled target plus
the layer-1 switch upweighting --- not more capacity.

\subsection{Optimizer and the closed-loop selection caveat}
AdamW, learning rate $3\!\times\!10^{-4}$, weight decay $10^{-2}$, gradient clip $1.0$,
batch $64$, $50$ epochs. We log, per epoch: per-modality validation losses, per-skill
action NLL, per-term backbone gradient norms and their cosine (to detect objectives
fighting over the trunk), the logic-gate precision/recall, and a \emph{switch
analytic} on the validation set --- at switch-decision positions, the skill head's
$P(\text{new})$, $P(\text{stay})$, and $\argmax$ accuracy. The offline validation loss
\emph{anti-correlates} with closed-loop success, so models are selected by closed-loop
rollout, not by held-out loss; the switch analytic ($\code{acc}:0\to0.85$ for the
decoupled run vs.\ $0\to0$ for the coupled one) predicts closed-loop behavior well.

% =========================================================================
\section{Inference}
\label{sec:inference}

At test time the transformer drives the BAMDP autoregressively, re-reading its own
skill head every step (Algorithm~\ref{alg:infer}, \code{rollout\_record}). There is no
discriminator and no logic gate; the only adaptation signal is the value token that
the BAMDP feeds back.

\begin{algorithm}[t]
\caption{Closed-loop rollout with on-line skill switching (\code{rollout\_record})}
\label{alg:infer}
\begin{algorithmic}[1]
\State history $\gets[\,]$;\ \ active $\gets \varnothing$
\While{episode not done}
  \State append state token $\st_t$ to history; \ $h\gets M(\text{history})$
  \State $\skill_t \gets \argmax M_{\mathrm{skill}}(h)$
  \Comment{the committed skill (argmax of the skill head)}
  \If{active $=\varnothing$}
     \State active $\gets \skill_t$; \ \textsc{SetActiveSkill}($\skill_t$)
  \ElsIf{$\skill_t \neq$ active}
     \State \textsc{SwitchSkill}($\skill_t$)
     \Comment{teleport to a fresh start (reset-on-switch), keep clock}
     \State $\st_t \gets \textsc{Observe}().\code{state}$; \ overwrite the last state
            token with the fresh start
     \State active $\gets \skill_t$; \ record a switch at $t$
  \EndIf
  \State append skill token $\skill_t$;\quad
         $a_t \gets \operatorname{clip}\big(\text{bin-center}\,\argmax\,
         M_{\mathrm{act}}(h),-1,1\big)$
  \State append value token $v^{\mathrm{obs}}_t$;\quad
         $o,\,\text{done}\gets\textsc{BamdpStep}(a_t)$;\quad $t\gets t{+}1$
\EndWhile
\end{algorithmic}
\end{algorithm}

\paragraph{Commitment and switching.}
Because the skill token is fed back into context, the model behaves as
\emph{sample-and-commit}: once a skill is in the history the skill head repeats it
(empirical stay probability $\approx0.993$) until the value stream stalls, at which
point the head flips to the good skill --- a switch. Under reset-on-switch the env
teleports to a fresh start, so the post-switch action must be computed from the
\emph{re-observed} fresh state; the implementation overwrites the just-appended state
token after the teleport so action and state are aligned. (Skipping this realignment
silently yields a policy that ``switches'' but never recovers, an instructive bug.)

\paragraph{What triggers a switch.}
Nothing in the loop references the goal or the latent. The switch is driven purely by
the exposed value $v^{\mathrm{obs}}_t$ flattening: in context, a sequence of value
tokens that climbs and then goes flat is the learned cue for ``abandon this skill.''
That is exactly the corruption signature \eqref{eq:corrupt} the model was trained on.

% =========================================================================
\section{Validation: transfer to a physical obstacle}
\label{sec:validation}

The policy is trained \emph{only} on synthetic value corruption. To check that it
learned a genuine stall$\to$switch reflex rather than memorizing the injection
statistics, we evaluate it with \emph{no injected failures} ($\theem=\mathbf 0$) in an
environment carrying a real obstacle: a circular \emph{freeze-on-contact} region
(\code{WallPoint2DGoalEnv}). The first time the point enters the region it is frozen
permanently (velocity zeroed, position held), so a skill whose path runs through it
gets physically stuck; its state stops progressing and $\Vnorm$ flatlines --- the
\emph{real} version of the synthetic failure. The obstacle is placed on one skill's
path while sparing the other (verified by rollout).

Across $60$ such walls the policy reaches the goal and switches to the unobstructed
skill in $60/60$ trials, with switches occurring at a mean step of $25.4$ (never in
the first $5$ steps) and the exposed value flat ($\Delta v\approx 0$ over the $5$ steps
before the switch) at a stalled level $\approx0.31\ll 1$. The behavior is symmetric:
forcing the policy to commit to the other skill and walling \emph{that} path yields the
reverse switch in $52/52$ trials. The synthetic-failure curriculum therefore transfers
to a physical stall it never saw in training.

% =========================================================================
\section{Summary}
\label{sec:summary}

\begin{enumerate}
\item A \emph{BAMDP} (\S\ref{sec:bamdp}) turns a fixed skill set into an adaptation
problem: a hidden per-skill failure rate $\theem$ drives a hazard process whose
absorbing failures corrupt an exposed, per-skill-normalized value so that a doomed run
is a value that stops climbing.
\item \emph{HG-DAgger} (\S\ref{sec:dagger}) collects on-policy episodes where the
transformer picks the skill, the BAMDP fails it, and an expert takes over with a good
skill after a short delay; the \emph{decoupled} skill target labels the switch a few
steps before the physical takeover.
\item A causal \emph{multimodal transformer} (\S\ref{sec:transformer}) is behavior-cloned
on state/skill/value tokens, predicting skill-from-state and action-from-skill, with
discretized actions, two layers of switch upweighting, and a logic-gate probe that
certifies the switch is decodable.
\item At \emph{inference} (\S\ref{sec:inference}) the model commits to a skill and
re-reads its skill head each step; a flat value stream flips the head to a better
skill, which under reset-on-switch retries from a fresh start --- and this transfers to
a real physical obstacle (\S\ref{sec:validation}).
\end{enumerate}

\end{document}
